% --------------------------------
% Unterschrift-Funktion
% --------------------------------

% Definiert einen neuen Befehl \signature zur Darstellung einer Unterschrift mit Ort und Datum.
% 
% Verwendung:
% \signature{Ort}
% 
% Dieser Befehl erzeugt eine Zeile mit dem angegebenen Ort, dem aktuellen Datum und einer Linie für die Unterschrift.
\newcommand{\signature}[2]{%
  \par\noindent
  \begin{tabular}{@{}l l@{}}
    % erste Spalte: feste Breite 6cm, Inhalt linksbündig
    \makebox[6cm][l]{#1, den \today} & \makebox[6.5cm]{\hrulefill}\\
    % zweite Zeile: keine erste Spalte, dann Unterschrift + Name
    & Unterschrift \textit{#2}
  \end{tabular}\par
}


% --------------------------------
% Definition der Seitenstile - Römische Seitenzahlen
% --------------------------------

% Definiert einen neuen Befehl \setupRomanPageStyle zur Konfiguration eines Seitenstils mit römischen Seitenzahlen.
%
% Verwendung:
% \setupRomanPageStyle
%
% Dieser Befehl stellt sicher, dass der Seitenstil auf fancyhdr umgestellt wird und die Seitenzahlen
% in römischen Ziffern angezeigt werden, ohne Kopfzeilen und mit der Seitenzahl unten rechts.
\newcommand{\setupRomanPageStyle}{
    \pagestyle{fancy} % Aktiviert das fancyhdr-Paket für Kopf- und Fußzeilen
    \fancyhf{} % Löscht alle aktuellen Einstellungen für Kopf- und Fußzeilen
    \fancyhead{} % Deaktiviert alle Kopfzeilenfelder
    \fancyfoot[R]{\thepage} % Setzt die Seitenzahl unten rechts in die Fußzeile
    \renewcommand{\headrulewidth}{0pt} % Entfernt die Trennlinie in der Kopfzeile
    \renewcommand{\footrulewidth}{0pt} % Entfernt die Trennlinie in der Fußzeile
    \pagenumbering{Roman} % Setzt die Seitennummerierung auf römische Ziffern
}

% --------------------------------
% Definition der Seitenstile - Arabische Seitenzahlen
% --------------------------------

% Definiert einen neuen Befehl \setupArabicPageStyle zur Konfiguration eines Seitenstils mit arabischen Seitenzahlen.
%
% Verwendung:
% \setupArabicPageStyle
%
% Dieser Befehl stellt sicher, dass der Seitenstil auf fancyhdr umgestellt wird und die Seitenzahlen
% in arabischen Ziffern angezeigt werden, mit dem Kapitel-/Abschnittstitel oben links und der Seitenzahl unten rechts.
\newcommand{\setupArabicPageStyle}{
    \pagestyle{fancy} % Aktiviert das fancyhdr-Paket für Kopf- und Fußzeilen
    \fancyhf{} % Löscht alle aktuellen Einstellungen für Kopf- und Fußzeilen
    \fancyhead[L]{\leftmark} % Setzt den Kapitel- oder Abschnittstitel oben links in die Kopfzeile
    \fancyfoot[R]{\thepage} % Setzt die Seitenzahl unten rechts in die Fußzeile
    \renewcommand{\headrulewidth}{1pt} % Legt die Dicke der Trennlinie in der Kopfzeile fest
    \renewcommand{\footrulewidth}{0pt} % Entfernt die Trennlinie in der Fußzeile
    \pagenumbering{arabic} % Setzt die Seitennummerierung auf arabische Ziffern
    % Setzt das Format für Abschnittstitel
}

% --------------------------------
% Flexible Tabelle
% --------------------------------

% Definiert einen neuen Befehl \flexTable zur Erstellung von flexiblen Tabellen mit benutzerdefinierten Einstellungen.
%
% Verwendung:
% \flexTable{Tabellenformat}{Tabellenkopf}{Tabelleninhalt}{Bildunterschrift}
%
% Dieser Befehl erzeugt eine Tabelle mit einer benutzerdefinierten Struktur, einem Tabellenkopf,
% dem Inhalt und einer Bildunterschrift. Die Tabelle wird innerhalb einer Gleitumgebung erstellt, 
% sodass sie auf der aktuellen Seite angezeigt wird.
\newcommand{\flexTable}[4]{
    \begin{table}[H] % Beginnt eine Tabelle in einer Gleitumgebung mit fester Position auf der Seite
    \centering % Zentriert die Tabelle auf der Seite
        \begin{tabularx}{\textwidth}{#1} % Erstellt eine Tabelle mit dem angegebenen Format
          \hline % Fügt eine horizontale Linie am Anfang der Tabelle ein
          #2 \\ % Fügt den Tabellenkopf ein
          \hline % Fügt eine horizontale Linie unterhalb des Tabellenkopfs ein
          #3 % Fügt den Tabelleninhalt ein
          \hline % Fügt eine horizontale Linie am Ende der Tabelle ein
        \end{tabularx}
    \caption{#4} % Fügt eine Bildunterschrift hinzu
    \label{tab:#4} % Setzt ein Label für die Tabelle, das für Verweise verwendet werden kann
    \end{table}
}

% --------------------------------
% Flexible Abbildung
% --------------------------------

% Definiert einen neuen Befehl \flexFigure zum Einfügen von Bildern mit benutzerdefinierten Einstellungen.
%
% Verwendung:
% \flexFigure{Breite}{Dateiname}{Bildunterschrift}
%
% Dieser Befehl fügt ein Bild in das Dokument ein und ermöglicht es, die Breite relativ zur Textbreite
% zu definieren. Eine Bildunterschrift und ein Label für Verweise werden ebenfalls hinzugefügt.
\newcommand{\flexFigure}[3]{
  \begin{figure}[H] % Beginnt eine Abbildung in einer Gleitumgebung mit fester Position auf der Seite
    \centering % Zentriert die Abbildung auf der Seite
    \includegraphics[width=#1\textwidth]{#2} % Fügt das Bild ein und skaliert es auf die angegebene Breite
    \caption{#3} % Fügt eine Bildunterschrift hinzu
    \label{fig:#2} % Setzt ein Label für die Abbildung, das für Verweise verwendet werden kann
  \end{figure}
}

% --------------------------------
% Flexible Quellcode-Einbettung
% --------------------------------

% Definiert einen neuen Befehl \flexCode zum Einfügen von Quellcode mit benutzerdefinierten Einstellungen.
%
% Verwendung:
% \flexCode{Sprache}{Bildunterschrift}{Dateiname}
%
% Dieser Befehl ermöglicht das Einfügen von Quellcode-Dateien in das Dokument, wobei die Sprache
% für das Syntax-Highlighting und eine Bildunterschrift angegeben werden. Ein Label für Verweise wird ebenfalls gesetzt.
\newcommand{\flexCode}[3]{
    \lstinputlisting
    [
        language=#1, % Legt die Programmiersprache für das Syntax-Highlighting fest
        caption={#2}, % Fügt eine Bildunterschrift für den Quellcode hinzu
        label={code:#3} % Setzt ein Label für den Quellcode, das für Verweise verwendet werden kann
    ]
    {#3} % Gibt den Dateinamen des Quellcodes an, der eingebunden werden soll
}